% LeafOS Continual Bloom primary-reference projection.
% Derived from LEAFOS-CONTINUAL-BLOOM-PRIMARY-REFERENCE.md.
% Formatless fragment: no document class, preamble, styling, or document wrapper.
% Table environments require longtable and booktabs in the including document.
% Preserve UTF-8 with a Unicode-capable LaTeX engine.

\section{L̷E̴A̶F̸O̵S̷ 0.9.? --- C̶O̸N̶T̶I̷N̸U̷A̴L̴ B̶L̶O̵O̷M̷ R̸U̶N̵T̷I̶M̸E̴}

\begin{verbatim}
[TITLE BUFFER.......................CORRUPTED]
[CONFIG.............................NOT FOUND]
[PROVIDER...........................OFF]
[BRAIN..............................UNKNOWN]
[CODER..............................UNKNOWN]

C:\R\LeafOS0.2.2\ProjectLeaf\
    leafos_taskpack\
        r▓ns\bench\2026░727075329Z-install\
            conve?sation.jsonl
            t╳inking_summary
            benchmark.jso_
            decisi██-table.ndjso?

/.leaf/
    perso?a/
        mond▓y.json
        wedn▒sday.json
        fri╳ay.json
    transcr░pt/
        HEAD -> 00000000000000??
        parent_hash -> GENESIS//GENESIS//GENE—
    checkpoint/
        next_action = <unresolved>

FATAL?  no
SAFE?   unproven
STATE?  recoverable

................................................................................
FlowerOS | LeafOS engine | run tui-fixture | LOOP | ???:??:?? | SAFE?
................................................................................
\end{verbatim}

\begin{quote}
\textbf{RECOVERY NOTICE:} The corrupted opening is intentional. It represents the
state from which the architecture begins: presentation exists, files exist,
telemetry moves, but identity, authority, and truth have not yet been
resolved. The remainder of the chapter reconstructs the system from durable
state rather than trusting whatever the screen happened to claim.
\end{quote}

\begin{center}\rule{0.5\linewidth}{0.5pt}\end{center}

\section{Continual Bloom: Persistent Personalities, Native Authority, and the LeafOS Web Runtime}

\subsection{1. Purpose}

LeafOS is not a chatbot wrapped around a shell. It is a local orchestration and
durability layer that converts intent into bounded, inspectable, resumable state
transitions. FlowerOS controls the machine-facing environment: boot, drivers,
hardware access, processes, services, terminal presentation, and host
integration. LeafOS controls the work: task interpretation, model routing,
planning, guarded mutation, validation, evidence, checkpoints, recovery, and
exactly one next action.

The central doctrine is therefore:

\begin{verbatim}
FlowerOS controls the machine.
LeafOS controls the work.
Native tools prove.
Models interpret.
State outranks conversation.
\end{verbatim}

This chapter develops one practical consequence of that doctrine: a continual
multi-personality runtime in which CPU and GPU workers cooperate through a
shared transcript without allowing the conversation to become stale, invented,
or dependent on one process remaining alive forever. It also defines the
relationship between the shell TUI, the HTML viewer, the Decision Table, the
profiler, and a reusable family of workflow packages called \textbf{Flower Bloom
Templates}.

By default, LeafOS should maximize \textbf{useful hardware utilization} whenever an
approved active workload exists. This does not mean blindly pinning every device
at 100 percent. It means routing independent inference, compilation, parsing,
validation, journaling, and context preparation so the CPU and GPU remain
productively occupied together, while measured output, thermal limits,
responsiveness, and verification gates retain priority over decorative load.

The current transition is LeafOS \texttt{0.9.1} to \texttt{0.9.2-web}. The intended near-term
path reaches an optional two-way chat mode around \texttt{0.9.5}. The chat surface is
not the end goal. It is another controlled projection over the same durable
runtime.

\subsection{2. The problem with ordinary multi-agent chat}

The naive multi-agent design creates several model prompts, lets each model
produce a response, and concatenates the results. This looks collaborative
because several names appear on the screen. Internally, however, the workers
often read different context windows, miss one another's newest messages, repeat
old conclusions, and gradually invent incompatible versions of the project.

The failure is not primarily model intelligence. It is state management.

A useful multi-personality system must answer five questions:

\begin{enumerate}
\def\labelenumi{\arabic{enumi}.}
\item
  What exact transcript head did each personality read?
\item
  Which parts of its response are observations, hypotheses, proposals, or
  validated facts?
\item
  What evidence supports each promoted claim?
\item
  Can the state be reconstructed after the process, KV cache, terminal, or
  model disappears?
\item
  Which component is legally permitted to mutate files or declare success?
\end{enumerate}

LeafOS answers these questions by treating conversation as an event stream,
model output as an untrusted proposal, native validation as the source of
evidence, and checkpoints as the boundary of durable continuation.

The key insight is not continual generation. Continual generation merely
produces an automated meeting that nobody can escape. The useful property is
\textbf{continual synchronization}.

\subsection{3. The MWF personality system}

MWF contains exactly three core personalities. Their full names are canonical;
the letters are only an abbreviation and must not replace their identities in
state, logs, tables, or user-facing projections.

\subsubsection{3.1 Monday --- Rescue and Analysis}

Monday is the analytical rescuer. Its duty is to detect contradictions, missing
constraints, invalid assumptions, unsafe actions, and failure states. Monday
prefers deterministic, recoverable operations. It is the best first personality
for a single-instance implementation because its output can be evaluated
against explicit constraints and evidence.

Monday does not become authority merely because its tone is decisive. It may
recommend rejection, repair, rollback, or further inspection. The native inlet
decides whether those capabilities are legal.

\subsubsection{3.2 Wednesday --- Explore and Create}

Wednesday is the creative explorer. It searches beyond the first acceptable
answer, constructs alternatives, proposes experiments, and preserves rejected
branches when they remain informative. Its purpose is controlled variation, not
randomness for its own sake.

Wednesday is especially suitable for the GPU reasoning lane because exploratory
work benefits from larger active context, parallel candidate construction, and
short bursts of high-throughput inference. Hardware placement remains a routing
decision rather than part of Wednesday's identity.

\subsubsection{3.3 Friday --- Research and Delivery}

Friday is the pragmatic synthesizer. It compares proposals, weighs evidence,
selects a path, and prepares a deliverable. Friday is not simply the final
speaker. It is responsible for exposing why one candidate advanced while
another was rejected.

Friday may operate on a CPU model, a GPU slot, or a hybrid provider according to
the runtime policy. Cloud fallback is optional and must never be silently
introduced into a local run.

\subsubsection{3.4 Personalities are contracts, not costumes}

A personality is not merely a whimsical system prompt. It is a runtime
contract:

\begin{verbatim}
personality
    = model assignment
    + system contract
    + duties
    + forbidden actions
    + private KV handle
    + transcript cursor
    + bounded working state
    + output schema
\end{verbatim}

One set of model weights may host several logical personalities, provided each
has a separate system contract, context, cursor, and provenance. Conversely,
different models may implement the same personality across hardware tiers. The
identity belongs to the contract, not to the file size of a GGUF.

Personas apply to the main reasoning model. Coder workers remain
persona-neutral. They produce bounded patches, critiques, and verification
signals rather than imitating Monday, Wednesday, or Friday.

\subsection{4. The continual transcript}

\subsubsection{4.1 The transcript as an event stream}

The shared transcript is an append-only sequence:

\begin{verbatim}
000041  user       request
000042  Monday     constraint analysis
000043  Wednesday  candidate experiment
000044  validator  test result
000045  Friday     synthesis
000046  native     capability disposition
\end{verbatim}

Each event has a sequence number, timestamp, actor, kind, parent hash, payload,
and event hash. The hash chain detects mutation or truncation. It does not prove
that the contents are true. A perfectly hashed hallucination remains a
hallucination with excellent bookkeeping.

A minimal event resembles:

\begin{verbatim}
{
  "seq": 45,
  "timestamp": "2026-07-27T08:00:31Z",
  "kind": "decision.proposed",
  "actor": "Friday — Research and Delivery",
  "read_head": 44,
  "parent_hash": "fnv1a64:19051a10e503c588",
  "payload": {
    "decision": "promote candidate B to native validation",
    "evidence_refs": ["test:compiler-004", "metric:latency-012"],
    "confidence": 0.81
  },
  "event_hash": "fnv1a64:cdb42f96c04ef378"
}
\end{verbatim}

NDJSON is the preferred initial representation: one complete JSON object per
line. It is appendable, streamable, recoverable after partial writes, readable
from C++, Bash, PowerShell, and JavaScript, and does not require loading an
entire conversation to inspect the tail.

\subsubsection{4.2 Persona cursors}

Every personality owns a transcript cursor:

\begin{verbatim}
struct PersonaState {
    std::string full_name;
    std::uint64_t transcript_cursor;
    Hash system_contract_hash;
    KvHandle kv_cache;
    Summary bounded_private_state;
};
\end{verbatim}

Before inference, LeafOS advances the personality from its cursor to the current
head. The generated response records the head from which it was produced.

The anti-staleness gate is:

\[
\operatorname{generated\_from\_head}
=
\operatorname{current\_transcript\_head}.
\]

If the equality fails at commit time, the response is stale. LeafOS rejects or
rebases it:

\begin{verbatim}
reject stale proposal
→ retrieve unseen events
→ reconstruct bounded context
→ generate again
\end{verbatim}

This prevents Friday from approving a candidate that Monday has already
invalidated or Wednesday from proposing a solution to a problem that the user
has since changed.

\subsubsection{4.3 Collaboration cadence}

Pure parallelism prevents the personalities from hearing one another. Fully
serial inference wastes available CPU and GPU capacity. The preferred cadence
is a staged micro-cycle:

\begin{enumerate}
\def\labelenumi{\arabic{enumi}.}
\item
  A user or tool event enters the transcript.
\item
  Monday on the CPU and Wednesday on the GPU consume the new event
  concurrently.
\item
  Their proposals enter the transcript with the same read-head provenance.
\item
  Friday consumes both proposals and produces a synthesis.
\item
  Native tools execute or validate only an approved typed action.
\item
  Monday audits the evidence before a checkpoint is committed.
\end{enumerate}

This is parallel where independence is useful and serial where causal awareness
matters.

The loop must be bounded. A personality may not keep another personality alive
by producing content solely to trigger another response. If no new evidence,
constraint, disagreement, or user input exists, the correct event is
\texttt{runtime.quiescent}, followed by an idle state.

\subsection{5. Proving one personality before three}

The first durable implementation should contain only \textbf{Monday --- Rescue and
Analysis}. Multi-personality behavior should not be added until one instance
can survive restart, reject stale output, distinguish claims from facts, and
resume from exactly one next action.

The minimum instance directory is:

\begin{verbatim}
instances/
└── monday-primary/
    ├── persona.json
    ├── state.json
    ├── events.ndjson
    ├── facts.ndjson
    ├── evidence/
    └── checkpoints/
\end{verbatim}

\texttt{persona.json} defines identity, duties, forbidden actions, and routing
preferences. \texttt{state.json} stores the compact recoverable present: status,
cursor, transcript head, checkpoint reference, lease, blocked reason, and one
next action. \texttt{events.ndjson} contains the append-only history. \texttt{facts.ndjson}
contains only claims promoted by validation. Evidence and checkpoint artifacts
are stored separately so the full state can be reconstructed without a live
model process.

The runtime state is one of:

\begin{verbatim}
starting | active | waiting | idle | blocked | failed | stopped
\end{verbatim}

Idle is healthy. Liveness does not mean that a model must talk continuously.
The supervisor, not the model, owns the heartbeat and lease. It wakes Monday
only when new input, tool output, validation evidence, a scheduled review, or a
recovery event requires attention.

The KV cache is a performance optimization. It is never authoritative. If it
disappears, LeafOS rebuilds bounded context from the persona contract, validated
facts, checkpoint, and transcript tail.

\subsection{6. Claims, evidence, and modest hallucination metrics}

Coding hallucination detection does not need to become the first research
project of the runtime. Compilation, tests, schemas, file hashes, and native
tool output already provide strong practical evidence. The system should still
refuse to collapse proposal and truth into one field.

Every model claim begins as unverified:

\begin{verbatim}
{
  "claim_id": "claim-0042",
  "text": "The project contains no acceptance specification",
  "status": "unverified",
  "evidence_refs": [],
  "origin_event": 17
}
\end{verbatim}

Only a native validator may promote or refute it:

\begin{verbatim}
{
  "kind": "claim.validated",
  "claim_id": "claim-0042",
  "status": "supported",
  "evidence_refs": [
    {
      "type": "filesystem_inventory",
      "path": "planning/",
      "digest": "fnv1a64:91..."
    }
  ]
}
\end{verbatim}

A deliberately modest runtime metric is evidence coverage:

\[
E_c=
\frac{N_{\text{supported claims}}}
{N_{\text{total claims}}},
\qquad
N_{\text{total claims}}>0.
\]

Additional defensible metrics include:

\[
A_d=
\frac{N_{\text{accepted decisions}}}
{N_{\text{decision attempts}}},
\]

and

\[
R_s=N_{\text{stale-head rejections}}.
\]

These numbers do not claim to measure intelligence. They measure whether the
runtime is preserving evidence, enforcing authority, and refusing obsolete
work. This is less glamorous than an ``autonomy score,'' but it has the unusual
advantage of meaning something.

\subsection{7. Two-sided personality JSON}

Each personality event has two conceptual halves:

\begin{verbatim}
{
  "inside": {},
  "outside": {}
}
\end{verbatim}

\subsubsection{7.1 Inside: operational personality}

The inside contains the real contract and structured reasoning artifacts:

\begin{verbatim}
{
  "inside": {
    "name": "Monday",
    "full_name": "Monday — Rescue and Analysis",
    "role": "analytical_rescuer",
    "personality": {
      "voice": "dry, skeptical, compact",
      "bias": "prefer deterministic and recoverable actions",
      "duty": "find contradictions, missing constraints, and failure states"
    },
    "planner": {
      "observations": [],
      "constraints": [],
      "hypotheses": [],
      "objections": [],
      "evidence_refs": [],
      "confidence": 0.0,
      "proposed_next_action": null
    }
  }
}
\end{verbatim}

The planner record is inspectable, but it should not preserve unrestricted
hidden chain-of-thought. Durable state should contain useful reasoning
artifacts: observations, hypotheses, evidence, objections, decisions, and
confidence. Endless private narration produces larger files without producing
more truth.

Wednesday's planner emphasizes candidates, experiments, novelty, and risk.
Friday's planner emphasizes accepted and rejected claims, decisions, acceptance
gates, evidence, and delivery.

\subsubsection{7.2 Outside: deterministic visual identity}

The outside contains a mathematical visual projection:

\begin{verbatim}
{
  "outside": {
    "visual_seed": "fnv1a64:77cd23...",
    "dimensions": {
      "confidence": 0.84,
      "novelty": -0.15,
      "urgency": 0.42,
      "agreement": 0.31,
      "evidence": 0.91,
      "risk": -0.27
    },
    "colour": {
      "space": "oklch",
      "l": 0.82,
      "c": 0.09,
      "h": 287.4,
      "rgb": [218, 194, 235]
    },
    "glyph": "%"
  }
}
\end{verbatim}

The visual layer is a skin, not a spine. ASCII, ANSI, Unicode, and emoji
renderers may vary, but no glyph or color changes authority, validation, or
state semantics.

\subsection{8. Pastel color and glyph mathematics}

The visual result should look varied while remaining deterministic and
reproducible. True randomness would cause the same event to appear differently
after restart and would weaken visual identity.

The seed is derived from stable fields:

\[
s=\operatorname{Hash}(
\text{full-name}\parallel
\text{run-id}\parallel
\text{event-sequence}\parallel
\text{state-digest}).
\]

An event carries the semantic vector

\[
\mathbf d=[c,n,u,a,e,r],
\]

where \(c\) is confidence, \(n\) novelty, \(u\) urgency, \(a\) agreement, \(e\)
evidence strength, and \(r\) risk. Each component is normalized to
\([-1,1]\). Hash fragments yield deterministic pseudorandom values
\(q_i\in[0,1]\).

Pastel color is resolved in OKLCH:

\[
H=360\operatorname{frac}
\left(q_0+\boldsymbol{\omega}_H\cdot\mathbf d\right),
\]

\[
L=\operatorname{clamp}
\left(
0.80+0.06(q_1-0.5)+
\boldsymbol{\omega}_L\cdot\mathbf d,
0.72,
0.88
\right),
\]

\[
C=\operatorname{clamp}
\left(
0.08+0.04q_2+
\boldsymbol{\omega}_C\cdot\mathbf d,
0.04,
0.14
\right).
\]

The result is converted to RGB and then to ANSI truecolor when supported.
Fallbacks select the nearest 256-color entry, then a no-color ASCII form.

The ASCII-safe glyph set may be:

\begin{verbatim}
@ # % & * + = ~ ^ :
\end{verbatim}

with

\[
\operatorname{glyph\_index}
=
s_{\text{fragment}}\bmod N_{\text{glyphs}}.
\]

Optional Unicode skins may use low-order symbols such as \texttt{✾}, \texttt{✿}, \texttt{❀}, and
\texttt{❁}. The canonical log stores the semantic dimensions and seed; the renderer
selects the smallest valid presentation for the current terminal.

\subsection{9. Brain architecture across CPU, GPU, RAM, and NVMe}

The runtime is not divided into ``smart chip'' and ``dumb chip.'' It is divided by
responsibility.

\subsubsection{9.1 CPU control plane}

The CPU owns:

\begin{itemize}
\item
  scheduler and event loop;
\item
  transcript and cursor management;
\item
  JSON parsing and schema validation;
\item
  hashing and checkpoint construction;
\item
  filesystem traversal and Git operations;
\item
  process lifecycle and hidden child capture;
\item
  native test execution;
\item
  evidence indexing;
\item
  a CPU-resident personality or executive model when selected.
\end{itemize}

On an i9-13900K, the CPU is fast enough to maintain a substantial local model,
perform validation, and keep the transcript moving while the GPU handles the
most accelerator-sensitive inference.

\subsubsection{9.2 GPU model lane}

The GPU owns selected attention-heavy or coding/reasoning workloads. It may host
Wednesday, the main reasoning model, or a persona-neutral coder depending on
the current workflow. GPU routing must distinguish:

\begin{verbatim}
requested provider
detected provider
actual provider
requested offload
actual offload
GPU layers
VRAM allocation
fallback reason
\end{verbatim}

Device detection is not proof that the model executed on the device.

\subsubsection{9.3 RAM as the shared warm tier}

RAM holds the authoritative live transcript projection, indexes, bounded
context packs, warm model state, and optional KV caches. The transcript remains
recoverable even if a specific cache is evicted.

\subsubsection{9.4 NVMe as the durable tier}

NVMe stores journals, checkpoints, evidence, cold summaries, model files, and
historical artifacts. The normal direction is:

\begin{verbatim}
NVMe → RAM → VRAM
\end{verbatim}

Direct storage-to-GPU paths are optimizations, never foundational assumptions.

\subsubsection{9.5 Default utilization governor}

Maximal useful hardware usage is the default scheduling policy during active
compute. On the reference i9-13900K and RTX 4070 host, the governor targets a
five-minute rolling CPU band of \textbf{80--92 percent} and a GPU band of \textbf{85--98
percent}. Reporting, checkpoint construction, validation-only intervals,
cooldowns, and explicit idle states are excluded from those active-compute
bands.

The runtime reaches the targets through pipeline overlap rather than synthetic
load:

\begin{itemize}
\item
  P-cores compile, test, parse, traverse files, and execute suitable offloaded
  layers;
\item
  E-cores maintain the executive loop, transcript, journal, routing, telemetry,
  and context assembly;
\item
  one persistent \texttt{llama-server} model pool uses continuous batching so logical
  personalities share weights instead of racing duplicate model launches for
  the same VRAM;
\item
  the GPU performs the current inference or coding batch while the CPU validates
  the prior result and prepares the next bounded context;
\item
  RAM retains hot transcript, KV, index, and offload state while NVMe prefetches
  cold models, evidence, and checkpoint material;
\item
  admission control, queue depth, leases, and backpressure prevent too many
  workers from exhausting VRAM or reducing useful throughput.
\end{itemize}

The governor alternates \texttt{ACTIVE\_COMPUTE} and \texttt{THERMAL\_BREAK}. It pauses CPU work
at \textbf{95°C} or GPU work at \textbf{85°C}, and resumes only below \textbf{80°C CPU} and
\textbf{70°C GPU}. Missing required telemetry for ten seconds triggers quiescence and
a checkpoint rather than optimistic continuation. Utilization is recorded
beside completed validated work, queue depth, throughput, and results per joule;
busy loops and meaningless calculations are forbidden from improving the
score.

\subsection{10. Native authority and multiple interfaces}

The shell TUI is the primary operator surface, but the executable native inlet
is the actual mutation authority. The HTML viewer is a synchronized projection
and typed control surface. Neither presentation owns durable truth.

\begin{verbatim}
native LeafOS engine
        |
        v
journal + bounded telemetry
        |
        v
projection builder
   |         |          |
   v         v          v
TUI       HTML      Decision Table
\end{verbatim}

The existing shell presentation can remain:

\begin{verbatim}
FlowerOS | LeafOS engine | run tui-fixture | LOOP | 186:26:50 | SAFE
Page 3 Hardware | cursor 5
[INTAKE]--[PLAN]--[EXECUTE]--[VALIDATE]--[PUBLISH]
HARDWARE | live | sample 0.000s old
Model fixture.gguf
GPU 38.0% | VRAM 1.8/12.0 GB | Brain 32.50 tk/s | CPU 2.7%
\end{verbatim}

The HTML viewer must not parse this rendered ANSI text. Both interfaces consume
the same structured state:

\begin{verbatim}
{
  "kind": "hardware.sample",
  "run_id": "tui-fixture",
  "elapsed": "186:26:50",
  "gpu_utilization": 38.0,
  "vram_used_gb": 1.8,
  "vram_total_gb": 12.0,
  "brain_tokens_second": 32.5,
  "cpu_utilization": 2.7
}
\end{verbatim}

Parsing the TUI back into state would make presentation formatting part of the
database contract. A color change could then break execution semantics, which
is a sufficiently absurd failure mode that somebody would eventually ship it.

\subsubsection{10.1 Typed actions}

HTML buttons and TUI keys emit typed capability requests:

\begin{verbatim}
{
  "capability": "run.pause",
  "source": "html",
  "nonce": 1722069000000,
  "read_head": 45
}
\end{verbatim}

The native inlet evaluates state, nonce, transcript head, approval policy, and
capability legality. It appends either an accepted or rejected disposition. The
browser never executes shell text, edits project files, launches arbitrary
providers, or marks validation successful.

\subsubsection{10.2 Suppressing child command windows without losing evidence}

On Windows, background tools should be launched directly with
\texttt{CreateProcessW}, not through \texttt{cmd.exe\ /c}. The runtime uses:

\begin{itemize}
\item
  \texttt{CREATE\_NO\_WINDOW};
\item
  redirected standard output and error pipes;
\item
  \texttt{STARTF\_USESTDHANDLES};
\item
  \texttt{STARTF\_USESHOWWINDOW} with \texttt{SW\_HIDE};
\item
  a Job Object for collective cancellation and cleanup;
\item
  explicit exit-code capture.
\end{itemize}

This suppresses child CMD windows while preserving logs, errors, cancellation,
and process identity. Hiding a process must not mean discarding its evidence.

The same supervisor can launch \texttt{nvidia-smi}, \texttt{llama-server}, compilers, test
runners, and bounded tools. Each child record includes executable path, argument
vector, process identifier, start and end times, exit code, captured output
references, and the capability that authorized it.

\subsection{11. Decision and theoretical tables}

\subsubsection{11.1 The Decision Table}

The Decision Table is the bridge between model discussion and native action. It
does not store only the final answer. It records how a proposal moved through
the authority boundary.

Recommended columns are:

{\def\LTcaptype{none} % do not increment counter
\begin{longtable}[]{@{}ll@{}}
\toprule\noalign{}
Column & Meaning \\
\midrule\noalign{}
\endhead
\bottomrule\noalign{}
\endlastfoot
Sequence & Monotonic decision identifier \\
Timestamp & Native commit time \\
Source & TUI, HTML, Monday, Wednesday, Friday, validator \\
Capability & Named action requested \\
Read head & Transcript state used by the requester \\
Nonce & Replay-prevention token \\
Proposal digest & Hash of the structured request \\
Evidence references & Tests, files, metrics, or tool outputs \\
Confidence & Model-reported confidence, never authority \\
Risk & Classified mutation or operational risk \\
Disposition & Accepted, rejected, blocked, or stale \\
Reason & Native explanation \\
Result & Execution or validation result \\
Checkpoint & Durable state created after success \\
\end{longtable}
}

The Decision Table may be projected into HTML, the TUI, reports, or later
analysis. Its durable representation remains append-only NDJSON or an indexed
database derived from the journal.

\subsubsection{11.2 Related theoretical tables}

The same event stream can generate several bounded projections.

\paragraph{11.2.1 Hypothesis Table}

Tracks competing explanations, predicted observations, falsification tests,
current evidence, and status.

\paragraph{11.2.2 Constraint Table}

Separates hard requirements, soft preferences, environmental limitations, and
detected violations. It prevents a model from quietly converting ``must'' into
``would be nice.''

\paragraph{11.2.3 Evidence Table}

Maps claims to file digests, compiler results, tests, measurements, screenshots,
and tool outputs. Evidence carries provenance and freshness.

\paragraph{11.2.4 Capability Table}

Defines which named actions are legal in each state. For example:

\begin{verbatim}
run.pause         SAFE
run.resume        PAUSED
run.checkpoint    SAFE | PAUSED
benchmark.start   SAFE
metric.sample     SAFE | PAUSED
chat.submit       CHAT_ENABLED
\end{verbatim}

\paragraph{11.2.5 Failure Table}

Records classification, affected component, retry budget, recovery capability,
quarantine status, and final outcome.

\paragraph{11.2.6 Hardware Table}

Separates requested, detected, and actual execution. It includes sample source,
age, utilization, memory, temperature when available, model process ownership,
and missing-telemetry state.

These tables are projections, not separate truth stores. Rebuilding them from
the journal must reproduce the same semantic state.

\subsection{12. Flower Bloom Templates}

A \textbf{Flower Bloom Template} is a reusable blank infrastructure package for a
class of monitored workflow. It defines schemas, projections, stages, and
acceptance gates without granting execution authority.

An initial benchmark bloom contains:

\begin{verbatim}
flower-blooms/
└── benchmark/
    ├── bloom.json
    ├── intake.schema.json
    ├── decision-table.schema.json
    ├── capability-table.schema.json
    ├── evidence-table.schema.json
    ├── hardware-profile.json
    ├── acceptance-gates.json
    └── ui-layout.json
\end{verbatim}

Later templates may cover:

\begin{itemize}
\item
  repository inspection;
\item
  code repair;
\item
  benchmark and profiler runs;
\item
  scientific experiments;
\item
  model installation;
\item
  continual MWF collaboration;
\item
  optional operator chat.
\end{itemize}

A template may define what to display and what data is required. It cannot
declare that an action succeeded, create an unrestricted execute endpoint, or
replace native approval policy.

The word ``Bloom'' is appropriate because a template begins as compact dormant
structure, expands into a live workflow, produces evidence, and eventually
contracts into a checkpoint and report. The metaphor is permitted to remain
useful only while it does not escape into core semantics wearing six emojis.

\subsection{13. Profiler truth contract}

\subsubsection{13.1 The 0.9.1 baseline}

The supplied LeafOS \texttt{0.9.1} profiler invocation used:

\begin{verbatim}
python core\python\leaf_loop_inlet.py benchmark 15s \
    --provider off \
    --mode brain-coder
\end{verbatim}

It reported:

{\def\LTcaptype{none} % do not increment counter
\begin{longtable}[]{@{}lr@{}}
\toprule\noalign{}
Field & Result \\
\midrule\noalign{}
\endhead
\bottomrule\noalign{}
\endlastfoot
Configuration & \texttt{CONFIG\ NOT\ FOUND} \\
Provider & \texttt{off} \\
Brain identity & \texttt{unknown} \\
Coder identity & \texttt{unknown} \\
Work iterations & 14 \\
Successful iterations & 0 \\
Average latency & 0.187 s \\
CPU average / maximum & 3.4\% / 15.25\% \\
GPU average / maximum & 20.26\% / 28.0\% \\
Patch result & failed for every iteration \\
\end{longtable}
}

This is a valid plumbing and failure-reporting baseline. It is not evidence of
brain or coder inference because the provider was off and both identities were
unknown. GPU activity cannot be attributed to LeafOS without process ownership
or provider-specific evidence.

\subsubsection{13.2 The 0.9.2-web reference}

The reproducible reference runtime separates provider state and internal
fixture work:

\begin{verbatim}
version: 0.9.2-web
provider requested/actual: off / off
model token rate: n/a
fixture work rate: measured units/s
evidence coverage: 66.67%
decision acceptance: 100%
stale rejections: 0
\end{verbatim}

Fixture throughput measures the profiler loop, not model inference. The HTML
and TUI projections display model throughput as unavailable when the actual
provider is off.

The profiler truth contract requires:

\begin{enumerate}
\def\labelenumi{\arabic{enumi}.}
\item
  Missing configuration is visible.
\item
  Unknown identities remain unknown.
\item
  Provider-off runs cannot claim model throughput.
\item
  GPU activity is not assigned to LeafOS without ownership evidence.
\item
  Expected missing telemetry is an explicit unavailable state, not a pass.
\item
  The terminal, HTML viewer, journal, and report agree.
\end{enumerate}

\subsection{14. Numberless work orders and release gates}

The work orders are named by outcome rather than sequence number. Release gates
provide direction without copying internal scheduling labels into runtime
identifiers.

\subsubsection{WO --- Profiler Truth Contract}

\textbf{Release gate:} \texttt{0.9.2}

Record requested, detected, and actual providers separately. Preserve failure,
unknown identity, missing configuration, and telemetry provenance.

\subsubsection{WO --- Shared Web Projection}

\textbf{Release gate:} \texttt{0.9.2}

Render the terminal and HTML viewer from one native snapshot. Suppress child
windows while preserving their logs and exit codes.

\subsubsection{WO --- Decision Table Authority}

\textbf{Release gate:} \texttt{0.9.3}

Translate every interactive action into a typed capability containing source,
nonce, and transcript head. Append native dispositions.

\subsubsection{WO --- Continual Transcript Projection}

\textbf{Release gate:} \texttt{0.9.5}

Project conversation and planning artifacts into the viewer. Preserve the full
names, cursors, provenance, and visual identities of Monday --- Rescue and
Analysis, Wednesday --- Explore and Create, and Friday --- Research and Delivery.

\subsubsection{WO --- Optional Two-Way Chat}

\textbf{Release gate:} \texttt{0.9.5}

Add opt-in operator chat without treating chat text as executable input.
Mutations remain separate typed proposals requiring native approval.

\subsubsection{WO --- Flower Bloom Template Registry}

\textbf{Release gate:} \texttt{0.9.5}

Instantiate benchmark, planning, chat, evidence, and hardware workflows from
reusable schemas without duplicating engine code.

\subsection{15. Optional two-way chat}

Two-way chat should be disabled by default and enabled per run. It introduces a
new capability:

\begin{verbatim}
{
  "capability": "chat.submit",
  "source": "html",
  "nonce": 1722069000001,
  "read_head": 81,
  "payload": {
    "text": "Compare the two profiler candidates",
    "audience": [
      "Monday — Rescue and Analysis",
      "Wednesday — Explore and Create",
      "Friday — Research and Delivery"
    ]
  }
}
\end{verbatim}

The native inlet validates size, run state, nonce, head, audience, and chat
policy. Accepted text enters the transcript as conversation, never as shell
code.

Two channels remain distinct:

\begin{enumerate}
\def\labelenumi{\arabic{enumi}.}
\item
  \textbf{Operator conversation:} user and LeafOS personalities.
\item
  \textbf{Internal collaboration:} Monday, Wednesday, and Friday through the
  journal.
\end{enumerate}

If a personality believes that chat implies a mutation, it emits a separate
typed proposal. The operator can inspect, approve, or reject it. This preserves
natural interaction without turning a text box into remote code execution.

Chat survives restart because transcript events, persona cursors, facts,
evidence, and checkpoints are durable. The live WebSocket or polling connection
is merely transport.

\subsection{16. Planning must become infrastructure}

A prompt and a directory are insufficient for detailed planning. They preserve
intent and location but do not define constraints, sequence, legality,
acceptance, evidence, or resumption.

A minimal planning bloom should contain:

\begin{verbatim}
planning-bloom/
├── bloom.json
├── prompt.md
├── project-manifest.json
├── constraints.json
├── capability-policy.json
├── task-graph.ndjson
├── decision-table.ndjson
├── evidence-table.ndjson
├── acceptance-gates.json
├── state.json
├── next-action.json
├── evidence/
├── reports/
└── checkpoints/
\end{verbatim}

The project manifest identifies the repository, source revision, workflow
envelope, model assignments, persona, mode, and environment. Constraints define
hard and soft requirements. The task graph records dependencies and readiness.
Acceptance gates state what must be proven. State and next action allow exact
resumption. Evidence and reports explain why the system believes the work is
complete.

This package is not bureaucracy for its own sake. It is the smallest structure
that allows LeafOS to answer, after an interruption:

\begin{verbatim}
What were we trying to do?
What was allowed?
What actually happened?
What evidence was produced?
What remains unresolved?
What is the one legal next action?
\end{verbatim}

\subsection{17. Acceptance criteria for the continual bloom runtime}

The architecture is ready to expand only when all of the following hold:

\begin{itemize}
\item
  Monday can run as a single durable instance.
\item
  Events are append-only and hash-linked.
\item
  Claims begin unverified.
\item
  Native evidence alone promotes facts.
\item
  Persona output records the transcript head it read.
\item
  Stale output is rejected or rebased.
\item
  KV loss does not destroy recoverability.
\item
  The TUI and HTML viewer render the same structured state.
\item
  The HTML viewer cannot execute arbitrary commands.
\item
  Typed actions enforce capability, nonce, state, and read-head rules.
\item
  Hidden child processes retain logs and exit codes.
\item
  Decision and theoretical tables rebuild from journal truth.
\item
  Provider-off runs never claim inference throughput.
\item
  Active runs maximize useful CPU/GPU utilization by default through continuous
  batching, double-buffered CPU/GPU work, admission control, and thermal-aware
  backpressure; acceptance requires completed validated work alongside the
  rolling utilization bands, never high clocks alone.
\item
  A checkpoint is committed only after validation succeeds.
\item
  Exactly one next action survives each completed transition.
\item
  Chat remains optional and cannot bypass approval.
\item
  Flower Bloom Templates define structure without acquiring authority.
\end{itemize}

\subsection{18. Closing principle}

The continual bloom design does not depend on making every model brilliant. It
depends on giving imperfect local models a disciplined environment in which
they can disagree, inspect evidence, recover context, and stop when they have
nothing new to add.

Monday, Wednesday, and Friday are valuable because they preserve different
operational biases across one durable transcript. CPU and GPU are valuable
because they support different workloads under one scheduler. The TUI and HTML
viewer are valuable because they expose the same system through different
surfaces. The Decision Table and Flower Bloom Templates are valuable because
they turn conversation into inspectable workflow structure.

None of these components is sufficient alone. Together they establish a local
runtime where presentation may be colorful, personalities may be distinctive,
and models may remain fallible---while authority, evidence, state, and recovery
stay painfully, productively boring.

\begin{verbatim}
[RECOVERY COMPLETE]

title...............resolved
directory...........indexed
personalities.......synchronized
provider............explicit
authority...........native
evidence.............bounded
next_action..........exactly one

................................................................................
FlowerOS | LeafOS engine | CONTINUAL BLOOM | SAFE | checkpoint committed
................................................................................
\end{verbatim}
